\documentclass{article}

\usepackage{pxLST}

\begin{document}

\section*{pxLST smoke + palette + alias regression}

% --- T1: Smoke load + environment coverage ---
\pxLSTsetup[dark]{pxJava}

\begin{CODE}{Smoke-CODE}
public class A { int x = 1; } // àèìòù áéíóú ç ñ ¿ ¡ · º ª
\end{CODE}

\begin{Code}[numbers=none]{Smoke-Code}
int add(int a, int b) { return a + b; }
\end{Code}

\begin{pxCode}{Smoke-pxCode}
public static void main(String[] args) {
    System.out.println("Hello P3CTeX!");
}
\end{pxCode}

\begin{pxCode*}{Smoke-pxCodeStar}
// Darker variant
int answer = 42;
\end{pxCode*}

\begin{miniCode}[.6\linewidth]
#!/bin/bash
echo "mini"
\end{miniCode}

\begin{CodeBox}{bx-one}{Title}{note}{Caption panel text.}
Funcio factorial(n: Integer): Integer
  Si n <= 0 Llavors
    Retornar 1
  Fi Si
Fi Funcio
\end{CodeBox}

\begin{pxCodeBox}{bx-px}{pxCodeBox Title}{pxnote}{pxCodeBox caption panel text.}
int value = 100;
\end{pxCodeBox}

% Force multiple pages (breakable) for CodeBox* and pxCodeBreak.
\begin{CodeBox*}{bx-break}{Long box}{note}{This must be breakable across pages.}
line 01
line 02
line 03
line 04
line 05
line 06
line 07
line 08
line 09
line 10
\end{CodeBox*}

\begin{pxCodeBreak}{bx-break2}{Long box 2}{note2}{Breakable alias.}
line A
line B
line C
line D
line E
\end{pxCodeBreak}

% --- T2: Palette switching yields correct colors ---
\pxLSTsetup[pastel]{pxJava}
\begin{pxCode}{Pastel-Palette}
String greeting = "pastel";
\end{pxCode}

{\pxLSTsetup[light]{pxYAML}%
 \begin{Code}{Scoped-light}
key: "value"
\end{Code}
}

\end{document}
